% ============================================================================
% model.tex
% A sequential-market model of granularity reforms (Forward -> Spot),
% built on Ito-Reguant (2016) and the residual-demand setup of the project
% proposal. Four arbitrage regimes (no / full / limited / strategic) per
% I&R Table 1; two bidder types (strategic residual monopolist, operational
% renewable aggregator); three states traversed along the auction-side
% path $(1,1)$ -> $(1,Q)$ -> $(Q,Q)$. Main body collects mechanisms,
% intuitions, and key formulas; complete derivations in appendix.
% ============================================================================

\documentclass[11pt]{article}

\usepackage[margin=2.4cm]{geometry}
\usepackage{graphicx}
\usepackage{amsmath, amssymb, amsthm}
\usepackage{booktabs}
\usepackage{array}
\usepackage{xcolor}
\usepackage{tikz}
\usetikzlibrary{positioning, arrows.meta, calc, shapes.geometric, decorations.pathreplacing}
\usepackage{hyperref}
\hypersetup{colorlinks=true, linkcolor=blue!50!black, citecolor=blue!50!black, urlcolor=blue!50!black,
            pdftitle={A sequential-market model of MTU granularity reforms},
            pdfauthor={Pablo Paramio Perez}}
\usepackage{microtype}
\usepackage[round,authoryear]{natbib}
\usepackage{enumitem}
\setlist[itemize]{leftmargin=*, itemsep=2pt, topsep=2pt}
\usepackage{titlesec}
\titlespacing*{\section}{0pt}{12pt}{4pt}
\titlespacing*{\paragraph}{0pt}{6pt}{4pt}
\renewcommand{\arraystretch}{1.1}
\linespread{1.06}

\newtheorem{proposition}{Proposition}

\title{\vspace{-1.5cm}\large A sequential-market model of granularity reforms\\[-2pt]\normalsize Quantity shifts, richer ladders and within-hour dispersion under full and strategic arbitrage\vspace{-0.4cm}}
\author{\normalsize Pablo Paramio P\'erez}
\date{\normalsize\today}

\begin{document}
\maketitle
\vspace{-0.8cm}

% =============================================================================
%  MAIN BODY  (target: 6 pages, fits as one section of the thesis)
% =============================================================================

\section{Theory: a sequential-market model of granularity reforms}\label{sec:theory}

\subsection{Setup: three sequential auctions, three agents}\label{sec:setup}

We extend \citet{ito_reguant_2016}'s sequential-markets framework on two dimensions --- an explicit operational aggregator and a granularity selector $Q^m \in \{1, Q\}$ at each stage --- and trace the Spanish backward reform path $(1,1) \to (1,Q) \to (Q,Q)$. The headline mechanism is that aggregator quantity shifts to the more granular market in proportion to within-hour forecast-error variance; equilibrium prices, bid ladders, cleared MWh, and within-hour wedge dispersion all adjust in response. Full derivations live in the appendix; the main body collects the comparative-statics output of each state along the path and the eight mechanisms that explain the empirical patterns of the companion paper.

A representative clock-hour contains $Q$ equal-length subperiods. Three sequential stages operate in order: a \textbf{forward} auction (Stage 1, $m=F$), a \textbf{spot} auction (Stage 2, $m=S$), and ex-post \textbf{balancing} settlement (Stage 3, $m=B$). Each stage has its own \emph{granularity} $Q^m \in \{1, Q\}$ --- it clears either once per hour or once per subperiod --- defining a $2^3 = 8$-state policy space. The main analysis fixes the auction-side path $(Q^{F}, Q^{S}) = (1,1) \to (1,Q) \to (Q,Q)$. The forward and spot markets map empirically to the electricity day-ahead and intraday auctions respectively; we use the more general terminology throughout. Figure~\ref{fig:model-structure} summarises.

\begin{figure}[h]
\centering
\hspace*{-2.6cm}%
\begin{tikzpicture}[
  node distance=0pt,
  stage/.style={draw, rounded corners=4pt, fill=blue!8, thick, minimum width=3.0cm, minimum height=1.6cm, align=center, font=\small},
  pathstate/.style={draw, rounded corners=3pt, fill=red!10, minimum width=1.2cm, minimum height=0.6cm, font=\small\bfseries, align=center},
  rowhead/.style={font=\small\bfseries, anchor=east, align=right},
  rolelbl/.style={font=\scriptsize, anchor=north east, align=right, gray!45!black, inner sep=1pt},
  collbl/.style={font=\scriptsize\itshape, anchor=south, align=center, gray!45!black},
  act/.style={font=\small\bfseries, anchor=center, align=center},
  arrtime/.style={->, >=Stealth, thick, gray!50!black},
  arrarb/.style={->, >=Stealth, ultra thick, orange!80!black},
  arrpath/.style={->, >=Stealth, thick, red!55!black},
  infolabel/.style={font=\footnotesize\itshape, gray!55!black, fill=white, inner sep=2pt},
]

% ===== Top panel: stages on a timeline =====
\node[stage] (F) at (0, 0)   {\textbf{Forward}\\[2pt] {\footnotesize Stage 1}\\[1pt] {\footnotesize $Q^{F}\!\in\!\{1,Q\}$}};
\node[stage] (S) at (4.5, 0) {\textbf{Spot}\\[2pt]     {\footnotesize Stage 2}\\[1pt] {\footnotesize $Q^{S}\!\in\!\{1,Q\}$}};
\node[stage] (B) at (9, 0)   {\textbf{Balancing}\\[2pt] {\footnotesize Stage 3}\\[1pt] {\footnotesize $Q^{B}\!\in\!\{1,Q\}$}};

\draw[arrtime] (F.east) -- node[above=1pt, infolabel]{$\mu_t$ realises} (S.west);
\draw[arrtime] (S.east) -- node[above=1pt, infolabel]{$\varepsilon_t$ realises} (B.west);

% ===== Middle panel: agent participation matrix =====
% Horizontal separator
\draw[gray!30, thick] (-5.6, -1.5) -- (10.6, -1.5);

% Column headers (under each stage column)
\node[collbl] at (0,   -1.7) {Forward};
\node[collbl] at (4.5, -1.7) {Spot};
\node[collbl] at (9,   -1.7) {Balancing};

% Row 1: Strategic bidder
\node[rowhead] at (-2.0, -2.3) {Strategic bidder};
\node[rolelbl] at (-2.0, -2.55) {residual monopolist; \\ produces \emph{without} shock};
\node[act, blue!60!black] at (0,   -2.3) {\large $\bullet$};
\node[act, blue!60!black] at (4.5, -2.3) {\large $\bullet$};

% Row 2: Aggregator
\node[rowhead] at (-2.0, -3.8) {Operational aggregator};
\node[rolelbl] at (-2.0, -4.05) {price-taker, $y_t=\bar y_t+\eta_t$; \\ bears imbalance $\gamma$};
\node[act, blue!60!black] at (0,   -3.8) {\large $\bullet$};
\node[act, blue!60!black] at (4.5, -3.8) {\large $\bullet$};
\node[act, blue!60!black] at (9,   -3.8) {\large $\bullet$};

% Row 3: Arbitrageur (with arrow connecting Forward and Spot)
\node[rowhead] at (-2.0, -5.3) {Arbitrageur};
\node[rolelbl] at (-2.0, -5.55) {position $s$ (or $\{s_t\}$); \\ regimes 1--4 of \citeauthor{ito_reguant_2016}};
\draw[arrarb] (-0.05, -5.3) -- (4.45, -5.3);
\node[act, orange!80!black] at (0,   -5.3) {\large $\bullet$};
\node[act, orange!80!black] at (4.5, -5.3) {\large $\bullet$};

% Horizontal separator (below)
\draw[gray!30, thick] (-5.6, -6.4) -- (10.6, -6.4);

% ===== Bottom panel: reform path =====
\node[rowhead, gray!55!black] at (-2.0, -7.0) {Reform path};
\node[pathstate] (s1) at (0,   -7.0) {$(1,1)$};
\node[pathstate] (s2) at (4.5, -7.0) {$(1,Q)$};
\node[pathstate] (s3) at (9,   -7.0) {$(Q,Q)$};
\draw[arrpath] (s1.east) -- node[above, font=\footnotesize, gray!45!black]{spot reform} (s2.west);
\draw[arrpath] (s2.east) -- node[above, font=\footnotesize, gray!45!black]{forward reform} (s3.west);

\end{tikzpicture}
\caption{\textbf{Sequential markets, granularity, and agents.} \textit{Top:} a forward market, a spot market, and a balancing settlement on a timeline; each has its own granularity selector $Q^m \in \{1, Q\}$ and information resolves between them. The forward market maps empirically to the electricity day-ahead auction, the spot market to the intraday auction. \textit{Middle:} which agents act at which stages (blue dots), and the arbitrageur position $s$ linking forward to spot (orange arrow). The strategic bidder produces without uncertainty and does not internalise the imbalance penalty; only the operational aggregator does, since it bears $\eta_t$. \textit{Bottom:} the auction-side reform path $(Q^F, Q^S) = (1,1) \to (1,Q) \to (Q,Q)$ traversed in the main analysis.}
\label{fig:model-structure}
\end{figure}

\paragraph{Bidders and primitives.} The \textbf{strategic bidder} is a residual monopolist with constant marginal cost $c$ and capacity $K$, which it can deliver \emph{exactly} (no production shock); it therefore does not internalise the imbalance penalty $\gamma$. It posts a two-tranche ladder per delivery product to discriminate against the within-hour demand shock $\mu_t$ (variance $\sigma_\mu^2$). The \textbf{operational aggregator} is a price-taker with stochastic production $y_t = \bar y_t + \eta_t$ (forecast-error variance $\sigma_\eta^2$); it allocates total hourly MWh $\bar Y_h$ across forward and spot and \emph{does} internalise the imbalance penalty, which scales with $\gamma(Q^B)$. The \textbf{arbitrageur} carries a net position $s$ from forward to spot, in one of the four regimes of \citeauthor{ito_reguant_2016} (\S\ref{sec:arbitrage}); arbitrage transmits price impact but not imbalance risk. The \citeauthor{ito_reguant_2016} residual-demand specification --- forward intercept $\bar A_t = \bar A + \mu_t$, spot demand inheriting the spread $p_1 - p_2$ with per-subperiod thinness $b_{2t}$, both net of $s$ and aggregator supply --- is stated in full in Appendix~\ref{app:setup}; we keep the primitives $(\sigma_\mu^2, \sigma_\eta^2, b_{mt})$ general.

\paragraph{Solution concept.} We solve the game by \textbf{backward induction}: at Stage~2 the strategic bidder takes the forward position as given and best-responds to spot residual demand; the arbitrageur takes the implied price spread and chooses $s$ within its regime; at Stage~1 the strategic bidder commits its forward ladder using rational expectations of the Stage-2 best responses. Full step-by-step backward induction in Appendix~\ref{app:state-60-60}--\ref{app:state-15-15}.

\paragraph{Renewable fringe and our relation to \citeauthor{ito_reguant_2016}.} \citeauthor{ito_reguant_2016}'s residual-demand interpretation already absorbs a renewable (or thermal) competitive fringe into the slope $b_m$: the strategic bidder serves $A - b_m p_m$, the implicit fringe serves the complement. Their setup is sufficient to deliver the strategic-bidder, arbitrage, and price-wedge mechanisms but is silent on \emph{why} the fringe's willingness to supply changes with granularity. We make the fringe explicit (the operational aggregator with $\sigma_\eta^2$ and an imbalance penalty $\gamma$) so the quantity-shift channel (M1) can be derived from first principles rather than absorbed into a granularity-dependent slope $b_m(Q^S, \sigma_\eta^2)$. We do not claim our setup is strictly preferable: it adds clarity on the quantity-shift mechanism at the cost of one extra agent. Reverting to the \citeauthor{ito_reguant_2016} interpretation --- $b_m$ parameterised on granularity and forecast variance --- delivers the same comparative statics.

\subsection{Arbitrage regimes: four benchmarks, one empirical}\label{sec:arbitrage}

Arbitrage closes the forward-spot price wedge only partially in practice. Real arbitrageurs have bounded capacity and internalise price impact, so the empirically relevant regime is limited arbitrage (\citeauthor{ito_reguant_2016} Result~3); the full and strategic-arbitrage benchmarks below are useful comparators but neither holds in data.

Following \citet[\S I.A and Table 1]{ito_reguant_2016}, arbitrage profit per matched product is $\pi^A(s) = s\,(p_1 - p_2)$. Four regimes pin $s$:
\begin{itemize}
\item[(1)] \emph{No arbitrage}: $s = 0$, positive forward premium $p_1 > p_2 > c$.
\item[(2)] \emph{Full / competitive arbitrage}: $s$ closes the spread, $p_1 = p_2$; the strategic bidder withholds more in response.
\item[(3)] \emph{Limited arbitrage}: a capacity ceiling $s \le K$ binds, sustaining a positive wedge whose size depends on $K$.
\item[(4)] \emph{Strategic arbitrage}: the arbitrageur internalises price impact and stops short of full, giving the closed-form benchmark
\begin{equation}
\boxed{\;p_1 - p_2 \;=\; \frac{p_1 - c}{4} \;>\; 0.\;}
\label{eq:strategic-arb}
\end{equation}
\end{itemize}
\textbf{Arbitrageurs have market power too} --- their participation and balance sheets are bounded. The empirically relevant regime in electricity is (3), limited arbitrage: full and strategic are useful benchmarks, but the persistent positive forward premium observed everywhere rules out (2), and per-product capacity binding rules out (4) at the granular state. Comparative-statics signs are the same across (1), (3), (4) (\citeauthor{ito_reguant_2016} Table 1).

\subsection{Pre-reform $(1,1)$: scarcity wedge, no within-hour dispersion}\label{sec:pre}

With both auctions hourly, no within-hour wedge dispersion is mechanically feasible --- each market has one product per clock-hour. Two features carry through every subsequent reform: the \citeauthor{ito_reguant_2016} scarcity wedge $p_1 > p_2 > c$ survives under any non-full-arbitrage regime, and both within-hour shocks ($\mu_t$ on demand, $\eta_t$ on aggregator production) load onto the balancing residual because neither hourly market can absorb them.

Both auctions hourly. The \citeauthor{ito_reguant_2016} scarcity wedge $p_1 > p_2 > c$ obtains; neither the hourly forward nor the hourly spot can absorb within-hour shocks, so the balancing discrepancy is $\Delta_t = \mu_t + \eta_t + \varepsilon_t$.

\subsection{Asymmetric state $(1,Q)$: four spot-side mechanisms}\label{sec:id15}

Spot granularity simultaneously activates four mechanisms: (M1) the aggregator shifts quantity to the granular market, in proportion to forecast-error variance; (M2) the granular-side price drops in response; (M3) within-hour wedge dispersion emerges from per-subperiod thinness asymmetry; and (M4) the strategic bidder's spot-side two-tranche ladder widens. The forward market is unchanged, so the balancing discrepancy still carries $\mu_t$.

The spot market becomes per-period; the forward market stays hourly.

\paragraph{(M1) The aggregator shifts quantity to the granular market.} Per-period spot scheduling matches per-subperiod expected production, collapsing within-hour imbalance to residual noise; the hourly forward cannot adjust per subperiod and incurs a mismatch cost. The aggregator therefore shifts $\Delta^*$ MWh from forward to spot, and \textbf{the shift is larger when within-hour forecast uncertainty is larger}:
\begin{equation}
\boxed{\;\partial \Delta^* / \partial \sigma_\eta^2 \;>\; 0.\;}
\label{eq:shift}
\end{equation}
Derivation in Appendix~\ref{app:aggregator}.

\paragraph{(M2) The granular-side price drops.} More aggregator supply on spot reduces the strategic bidder's spot residual demand. The equilibrium spot price drops on average by an amount of order $\Delta^*/b_2$; the forward price drifts up modestly as the strategic bidder serves a larger forward share. The net effect across both markets is positive-sum.

\paragraph{(M3) Block parity reveals within-hour spot dispersion.} The hourly forward can be arbitraged only against the \emph{average} of spot per-period prices --- under full arbitrage, $p_1 = (1/Q)\sum_t p_{2t}$ (block parity); under strategic arbitrage, the same average condition plus the \citeauthor{ito_reguant_2016} wedge $(p_1-c)/4$. This leaves room for within-hour spot dispersion whenever per-subperiod thinness $b_{2t}$ varies across products: thinner products carry higher clearing prices for the same aggregate block. The wedge SD inherits this thinness asymmetry directly, attenuated by half under strategic vs full arbitrage (Appendix~\ref{app:state-60-15}).

\paragraph{(M4) Spot-side bid ladders widen.} Per-period clearing rewards two-tranche discrimination across products: the strategic bidder posts a low tranche for low-$\mu_t$ states and a high tranche for high-$\mu_t$ states. The optimal pre-realisation spread is
\begin{equation}
\boxed{\;\Delta p^*_{2t} \;=\; \frac{\Delta_\mu}{2\, b_{2t}},\;}
\label{eq:tranche-id}
\end{equation}
strictly larger than the hourly counterpart (Appendix~\ref{app:tranche}). Empirical measures of within-product bid-shape dispersion (e.g.\ MW-weighted price SDs, effective tranche counts) are motivated by this widening but are not derived from the two-tranche reduction directly; we use them as empirical proxies in the companion paper, not as model objects.

Balancing discrepancy unchanged at $\Delta_t = \mu_t + \varepsilon_t$: the forward is still hourly so $\mu_t$ still loads onto balancing.

\subsection{Symmetric state $(Q,Q)$: limited per-product arbitrage}\label{sec:da15}

When the forward also becomes per-period, six mechanisms operate; the empirically critical one is per-product limited arbitrage. Per-product arbitrage capacity binds at $K_t \approx K/Q$, so the within-hour wedge SD does not collapse to the strategic-arbitrage benchmark but stabilises at $\sigma_\mu/(4 b_1)$ plus thinness asymmetry --- exactly twice the unlimited strategic benchmark and consistent with the empirical asymmetric-persistence finding. The Cournot scale-up of per-clock-hour cleared MWh and the cross-market option-value substitution of bid shapes complete the symmetric-state predictions.

The forward market also becomes per-period.

\paragraph{(i) Aggregator allocation equalises.} The forward-vs-spot hedging differential closes once both markets accept per-subperiod schedules.

\paragraph{(ii) Per-product arbitrage --- and arbitrageur market power.} The full-arbitrage benchmark gives $p_{1t} = p_{2t}$ per product (wedge SD zero); the strategic-arbitrage benchmark gives the \citeauthor{ito_reguant_2016} wedge \eqref{eq:strategic-arb} per product, collapsing the wedge SD to $\sigma_\mu/(8\, b_1)$. \emph{Both benchmarks rely on the arbitrageur being able to build any per-product position, and both fail empirically.} Arbitrage capacity $K$ is bounded; with $Q$ matched products per clock-hour, per-product capacity is $K_t \approx K/Q$, which binds tightly. The symmetric state is therefore in the \textbf{limited-arbitrage regime} (\citeauthor{ito_reguant_2016} Result 3) applied per product: the per-product wedge becomes $(p_{1t}-c)/2$ rather than $(p_{1t}-c)/4$, doubling the wedge SD to $\sigma_\mu/(4\,b_1)$ (Appendix~\ref{app:15-15:limited}); if per-period thinness $b_{2t}$ also varies across products this contributes an additional thinness-asymmetry component. The wedge SD therefore does \emph{not} collapse from the asymmetric-state magnitude. \emph{This is the model-internal structural account of empirical asymmetric persistence: arbitrageur market power generates the non-collapse. Residual dispersion the model does not pin down may reflect additional channels (e.g.\ behavioural inertia or renewable-penetration trends not absorbed by the covariates); the structural and residual accounts are complementary, not exclusive.}

\paragraph{(iii) Forward-side bid ladders widen.} The same mechanism as (M4) on the forward side: optimal spread $\Delta p^*_{1t} = \Delta_\mu/(2 b_1)$.

\paragraph{(iv) Cournot scale-up from per-period thinness.} Per-period clearing thins the per-product residual demand: each per-subperiod product has fewer competing units than the hourly aggregate, so $b_{1t} < b_1$ on average. Cournot quantity is decreasing in elasticity, so per-product cleared MWh exceeds $1/Q$ of the hourly value and total hourly cleared MWh rises. The condition is $\sum_t b_{1t} < Q\, b_1$ (Appendix~\ref{app:15-15:scaleup}); the scale-up concentrates in subperiods where per-period thinness drops most relative to the hourly aggregate. No external withholding mechanism is needed --- this is the standard Cournot quantity comparative static.

\paragraph{(v) Cross-market bid-shape substitution via Stage-1 option value.} Stage-1 forward bidding is optimised by backward induction against expected Stage-2 outcomes. Per-period spot makes Stage-2 profit convex in $\mu_t$ (per-period Cournot gives $\partial^2 \pi_{2t}/\partial \mu_t^2 = 1/(2 b_{2t})$), so Stage-2 expected profit acquires a positive variance correction $\sigma_\mu^2 / (4 b_{2t})$. The Stage-1 bidder responds by widening the forward ladder to capture both sides of the variation it anticipates downstream. Symmetrically, once the forward is per-period the within-hour discrimination is absorbed at Stage 1 directly and the residual spot ladder tightens. (Derivation in Appendix~\ref{app:15-15:option}.)

\paragraph{(vi) Balancing efficiency.} The per-period forward schedule absorbs $\mu_t$, so only the residual post-Stage-2 noise loads onto the balancing reference: $\Delta_t = \varepsilon_t$.

\subsection{Comparative statics across the reform path}\label{sec:speaks}

The model delivers seven comparative-statics predictions across the three states $(1,1) \to (1,Q) \to (Q,Q)$, summarised in Table~\ref{tab:comparative}. The wedge-SD row reports the empirically relevant limited-arbitrage benchmark; the full and strategic benchmarks are derived in eq.~\eqref{eq:wedge-sd-asym-full-strat} for comparison.

\begin{table}[h]
\centering
\footnotesize
\caption{Comparative statics across the reform path under limited per-product arbitrage (empirical regime). Imbalance row reports $E[\Delta_t^2]$ for the quadratic-cost penalty (\S\ref{sec:welfare}); under uniform $\mu_t \sim U[-\Delta_\mu, \Delta_\mu]$ and $\eta_t \sim U[-\Delta_\eta, \Delta_\eta]$, $E[\mu_t^2] = \Delta_\mu^2/3$ and $E[\eta_t^2] = \Delta_\eta^2/3$.}
\label{tab:comparative}
\begin{tabular}{p{4.6cm} c c c}
\toprule
Outcome & $(1,1)$ & $(1,Q)$ & $(Q,Q)$ \\
\midrule
Aggregator fwd$\to$spot shift
 & $0$
 & $\Delta^*(\sigma_\eta^2) > 0$
 & equalised \\
Spot ladder spread $\Delta p^*_{2t}$
 & ---
 & $\dfrac{\Delta_\mu}{2 b_{2t}}$
 & $\dfrac{\Delta_\mu}{2 b_{2t}}$ \\
Forward ladder spread $\Delta p^*_{1t}$
 & ---
 & via Stage-1 option
 & $\dfrac{\Delta_\mu}{2 b_1}$ \\
Wedge SD (limited arb.)
 & $0$
 & $\dfrac{K|b_{22}-b_{21}|}{8 b_{21} b_{22}}$ ($Q\!=\!2$)
 & $\dfrac{\sigma_\mu}{4 b_1}$ + thin. asym. \\
Cleared MWh / clock-hour
 & baseline
 & forward unchanged
 & $+\dfrac{c(Q b_1 - \sum_t b_{1t})}{2}$ \\
Balancing $\Delta_t$
 & $\mu_t+\eta_t+\varepsilon_t$
 & $\mu_t+\varepsilon_t$
 & $\varepsilon_t$ \\
Imbalance cost $C^I/(\gamma Q)$, quadratic
 & $(\Delta_\mu^2+\Delta_\eta^2)/3$
 & $\Delta_\mu^2/3$
 & $0$ \\
\bottomrule
\end{tabular}
\end{table}

\paragraph{Arbitrage regime is empirical.} Real arbitrageurs have market power and bounded balance sheets: the persistent forward premium observed across electricity markets rules out full arbitrage, and the per-product capacity constraint binds at $(Q,Q)$. The relevant \citeauthor{ito_reguant_2016} column is Result~3 (limited), not Result~4 (strategic, unlimited); per-product limited arbitrage at $(Q,Q)$ delivers the model's account of the empirical asymmetric-persistence of wedge dispersion (\S\ref{sec:da15}(ii)).

\paragraph{Predictions to test empirically.} The model delivers seven sign predictions across the reform path; the empirical section of the thesis verifies them in turn:
(P1) quantity shifts from forward to spot at the spot reform, larger in clock-hours with bigger forecast uncertainty;
(P2) the granular-side equilibrium price drops in response;
(P3) within-hour wedge dispersion emerges from per-subperiod thinness asymmetry (largest in ramp hours);
(P4) the strategic bidder's two-tranche ladders widen on the granular side;
(P5) total cleared MWh per clock-hour rises at the symmetric state;
(P6) within-hour wedge SD does not collapse at $(Q,Q)$ --- the empirical asymmetric-persistence finding;
(P7) production-unit arbitrageurs exist and are capacity-constrained (we test this directly by checking whether production plants submit both buy and sell offers in OMIE data).

\subsection{Welfare across reforms and arbitrage regimes}\label{sec:welfare}

Under uniform shocks $\mu_t \sim U[-\Delta_\mu, \Delta_\mu]$ and $\eta_t \sim U[-\Delta_\eta, \Delta_\eta]$, the reform's welfare effect decomposes into three positive channels and one distributional one. The dominant gain is the imbalance saving from finer scheduling; a secondary gain is the Cournot scale-up at $(Q,Q)$; a third per-product gain comes from limited arbitrage replacing the strategic-arbitrage benchmark, where \citet{ito_reguant_2016} Result~2 implies less arbitrage activity reduces strategic withholding and raises output. The richer two-tranche ladder is mostly distributional. The empirical asymmetric-persistence of within-hour wedge SD (\S\ref{sec:da15}(ii)) characterises the limited regime; it is not a welfare cost.

\paragraph{Components.} Per clock-hour, $W = CS + \pi^{SB} + \pi^{Agg,\text{net}} + \pi^{Arb} - C^I_{\text{sys}}$. We model the imbalance penalty as quadratic in the deviation ($\gamma \mathcal{I}^2$ per subperiod); $\pi^{Agg,\text{net}}$ is the aggregator's profit net of its own quadratic penalty on $\eta_t$, and $C^I_{\text{sys}}$ is the system-side cost from $\mu_t$. The two channels are disjoint by construction --- $\eta_t$ only enters $\pi^{Agg,\text{net}}$, $\mu_t$ only enters $C^I_{\text{sys}}$.

\paragraph{Gain: imbalance saving falls along the reform path.} Spot granularity absorbs $\eta_t$ via aggregator per-subperiod hedging; forward granularity absorbs $\mu_t$ via per-product forward clearing. Under uniform shocks ($E[\mu_t^2] = \Delta_\mu^2/3$, $E[\eta_t^2] = \Delta_\eta^2/3$) and quadratic penalty:
\begin{center}
\footnotesize
\begin{tabular}{l c c c}
\toprule
& $(1,1)$ & $(1,Q)$ & $(Q,Q)$ \\
\midrule
Aggregator-internalised cost (from $\eta_t$ when spot hourly) & $\gamma Q \Delta_\eta^2/3$ & $0$ & $0$ \\
System cost $C^I_{\text{sys}}$ (from $\mu_t$ when forward hourly) & $\gamma Q \Delta_\mu^2/3$ & $\gamma Q \Delta_\mu^2/3$ & $0$ \\
\textbf{Total imbalance cost} & $\boldsymbol{\gamma Q(\Delta_\mu^2+\Delta_\eta^2)/3}$ & $\boldsymbol{\gamma Q \Delta_\mu^2/3}$ & $\boldsymbol{0}$ \\
\bottomrule
\end{tabular}
\end{center}
\paragraph{Gain: Cournot scale-up at $(Q,Q)$.} Per-period elasticities are smaller than the hourly aggregate ($\sum_t b_{1t} < Q b_1$), so $(c/2)(Q b_1 - \sum_t b_{1t})$ additional MWh clear at the pre-reform mark-up, delivering a welfare gain $\Delta W^{\text{scale}} = c(\bar A - b_1 c - S^{op}_1)(Q b_1 - \sum_t b_{1t})/(4 b_1)$.

\paragraph{Gain: less strategic withholding under limited per-product arbitrage.} The spot reform stretches per-product arbitrage capacity to $K/Q$, moving the regime from strategic to limited per product. Per \citet{ito_reguant_2016} Result~2, less arbitrage means less strategic withholding: the strategic bidder's spot markup falls from $3(p_1-c)/4$ (strategic) to $(p_1-c)/2$ (limited), and the per-product deadweight loss falls by $\Delta W^{\text{lim-arb}} = 5 b_2(p_1-c)^2/32 > 0$. The empirical asymmetric-persistence of wedge SD (\S\ref{sec:da15}(ii)) is the equilibrium signature of this regime change, not a welfare cost.

\paragraph{Net effect.} All channels positive: $\Delta W^{\text{spot}} = \gamma Q \Delta_\eta^2/3 + (\text{per-product lim-arb gain})$; $\Delta W^{\text{forward}} = \gamma Q \Delta_\mu^2/3 + \Delta W^{\text{scale}}$. The bid-ladder rent $\Delta\pi^{2T} = \Delta_\mu[4(\bar A - bc) + \Delta_\mu]/(16 b)$ (\S\ref{app:tranche}) is mostly a CS-to-$\pi^{SB}$ transfer. The reform is welfare-improving at each step.

{\footnotesize
\bibliographystyle{plainnat}
\bibliography{references}
}

\newpage

% =============================================================================
%  APPENDIX --- FULL ALGEBRA
% =============================================================================

\appendix
\section{Setup details and information arrival}\label{app:setup}

This appendix records the timing of information arrival, the \citeauthor{ito_reguant_2016} residual-demand specification used throughout, and how that specification specialises when one or both markets are hourly.

\paragraph{Timing.} Stage~1 (DA): only $\bar A$ and the distribution of $\mu_t$ are known to the strategic bidder; the aggregator knows $\bar Y_h$, $\{\bar y_t\}$ and the distribution of $\eta_t$. Stage~2 (ID): $\mu_t$ has realised and is common knowledge to first order. Stage~3 (balancing): $\eta_t$ realises and settlement occurs at the balancing reference.

\paragraph{Residual demand (\citet{ito_reguant_2016} form).} The strategic residual monopolist with constant marginal cost $c$ serves what the operational aggregator does not cover, net of any arbitrage position $s_t$ carried from forward to spot. Per delivery product:
\begin{equation}
q_{1t} \;=\; \bar A_t \;-\; b_{1t}\, p_{1t} \;-\; s_t \;-\; S^{op}_{1t},
\qquad
q_{2t} \;=\; b_{2t}\,(p_1^{(t)} - p_{2t}) \;+\; s_t \;-\; S^{op}_{2t},
\label{eq:rd}
\end{equation}
where $\bar A_t = \bar A + \mu_t$ is the per-subperiod demand intercept, $p_1^{(t)} = p_1$ when DA is hourly and $p_{1t}$ when DA is per-period, $S^{op}_{mt}$ is the aggregator's offered MWh in market $m \in \{1, 2\}$ at subperiod $t$ (uniform $S^{op}_m / Q$ when market $m$ is hourly), and $b_{mt}$ is market thinness (with per-subperiod asymmetry allowed). The within-hour demand shock satisfies $\mathbb{E}[\mu_t] = 0$ and $\mathrm{Var}(\mu_t) = \sigma_\mu^2$; the aggregator's per-subperiod realisation is $y_t = \bar y_t + \eta_t$ with forecast-error variance $\sigma_\eta^2$.

\paragraph{Hourly residual demand by specialisation.} The general residual demand \eqref{eq:rd} specialises by setting $\bar A_t \equiv \bar A$, $b_{mt} \equiv b_m$ and $S^{op}_{mt} \equiv S^{op}_m / Q$ when market $m$ is hourly. In $(1,1)$ both reduce to one product apiece; in $(1,Q)$ only the spot demand keeps the per-period index.

\paragraph{Empirical scope.} The two-stage analysis here covers the forward-spot auction pair. The balancing market enters only as the ex-post settlement reference; we collapse it into the discrepancy $\Delta_t$ used in \S\ref{app:balancing}.

\section{Operational aggregator: full scheduling problem}\label{app:aggregator}

We solve the aggregator's full scheduling problem in four steps: the problem statement, the closed-form equilibrium when both markets are hourly, the equilibrium under spot-per-period scheduling, and the comparative static $\partial \Delta^*/\partial \sigma_\eta^2 > 0$ that closes \eqref{eq:shift}.

\subsection{Problem statement}\label{app:aggregator:problem}
The aggregator chooses $(S^{op}_{1t}, S^{op}_{2t})_{t=1}^{Q}$ to solve
\begin{equation}
\max_{(S^{op}_{1t},\, S^{op}_{2t})_t}\;
\mathbb{E}\!\left[\sum_t \big(p_1^{(t)} S^{op}_{1t} + p_{2t} S^{op}_{2t}\big) \;-\; \gamma\,\mathcal{I}\right]
\quad\text{s.t.}\quad \sum_t (S^{op}_{1t} + S^{op}_{2t}) = \bar Y_h,
\label{eq:agg-problem}
\end{equation}
with quadratic imbalance penalty $\mathcal{I} = \sum_t (y_t - S^{op}_{1t} - S^{op}_{2t})^2$ under per-period settlement and uniform-across-$t$ allocation in any market that is hourly.

\subsection{Equilibrium under hourly schedules ($(1,1)$ and $(1,Q)$ with hourly aggregator)}\label{app:aggregator:hourly}
With both markets hourly the aggregator chooses scalar $(S^{op}_1, S^{op}_2)$ with $S^{op}_1 + S^{op}_2 = \bar Y_h$. The per-subperiod implementation is uniform: $S^{op}_{1t} = S^{op}_1/Q$ and $S^{op}_{2t} = S^{op}_2/Q$. Per-subperiod imbalance under per-period settlement is
\[
\mathcal{I} \;=\; \sum_t \big(y_t - (S^{op}_1 + S^{op}_2)/Q\big)^2 \;=\; \sum_t \big((\bar y_t - \bar Y_h/Q) + \eta_t\big)^2,
\]
which depends on $(\bar y_t - \bar Y_h/Q)$ and $\eta_t$, not on the forward-vs-spot split. The aggregator's allocation is therefore driven purely by the price differential: it commits at forward if $p_1 > p_2$ and indifferent if $p_1 = p_2$. We normalise $S^{op}_1 = S^{op}_2 = \bar Y_h / 2$ as a tie-break.

\subsection{Equilibrium under DA-hourly + ID-granular (state $(1,Q)$)}\label{app:aggregator:60-15}

\paragraph{Notation for forecast resolution.} The within-hour heterogeneity that loads onto the aggregator's imbalance has two components: a structural per-subperiod expected production variation $(\bar y_t - \bar Y_h/Q)$, known already at Stage~1, and a forecast update $\eta_t$ resolved between Stages~1 and~2 (i.e.\ between the forward close and the spot close). $\eta_t$ has variance $\sigma_\eta^2$ and is the empirically relevant ``forecast-error shock'': bigger $\sigma_\eta^2$ means more within-hour production information that materialises after the forward allocation is fixed.

\paragraph{Optimal per-subperiod spot schedule.} The aggregator chooses per-subperiod $\{S^{op}_{2t}\}_t$ at spot using its Stage-2 forecast $\bar y_t^{(2)} \equiv \bar y_t + \eta_t$ (the updated expected production at the time of the spot close). Setting
\begin{equation}
S^{op}_{2t}{}^{*} \;=\; \bar y_t^{(2)} - \frac{S^{op}_1}{Q} \;=\; \bar y_t + \eta_t - \frac{S^{op}_1}{Q},
\label{eq:agg-perperiod-id}
\end{equation}
matches the Stage-2 forecast net of the uniformly-allocated forward commitment. The realised per-subperiod imbalance under \eqref{eq:agg-perperiod-id} reduces to the residual post-Stage-2 noise $\varepsilon_t$ (independent of $\eta_t$, with small variance), so the expected aggregate imbalance under per-period scheduling is
\[
\mathbb{E}\!\left[\,\mathcal{I}_{\text{granular}}\,\right] \;=\; Q\cdot\mathbb{E}\!\left[\,|\varepsilon_t|\,\right].
\]
Under the hourly schedule (\S\ref{app:aggregator:hourly}) the aggregator must commit at Stage~1 using the Stage-1 forecast $\bar y_t^{(1)} = \bar y_t$; the Stage-1-to-Stage-2 update $\eta_t$ then loads onto the imbalance:
\[
\mathbb{E}\!\left[\,\mathcal{I}_{\text{hourly}}\,\right] \;=\; \sum_{t=1}^{Q} \mathbb{E}\!\left[\,\big|(\bar y_t - \bar Y_h/Q) + \eta_t + \varepsilon_t\big|\,\right].
\]

The aggregator's net gain from per-period spot scheduling is
\begin{equation}
G(\sigma_\eta^2) \;\equiv\; \gamma\!\sum_t\!\Big[\mathbb{E}\big|(\bar y_t - \bar Y_h/Q) + \eta_t + \varepsilon_t\big| \;-\; \mathbb{E}\big|\varepsilon_t\big|\Big],
\label{eq:agg-gain}
\end{equation}
strictly increasing in $\sigma_\eta^2$ (the avoided Stage-1-to-Stage-2 forecast shock) and in $\sigma_a^2 \equiv \mathrm{Var}(\bar y_t - \bar Y_h/Q)$ (the structural within-hour variation). Both components map empirically to within-hour heterogeneity: clock-hours where production swings within the hour carry larger $\sigma_\eta^2$ and larger $\sigma_a^2$.

\subsection{Quantity shift $\Delta^*(\sigma_\eta^2)$ and its comparative static}\label{app:aggregator:shift}
The aggregator's marginal trade-off in $S^{op}_1$ is the forward-vs-spot price differential against the marginal change in $G$. The first-order condition gives an implicit $\Delta^*(\sigma_\eta^2)$:
\begin{equation}
(p_1 - p_2) \;=\; \gamma \cdot \frac{\partial}{\partial S^{op}_1}\!\bigg[\sum_t \mathbb{E}\big|(\bar y_t - (\bar Y_h - \Delta^*)/Q) + \eta_t\big|\bigg]\bigg|_{\Delta^*}.
\label{eq:agg-foc}
\end{equation}
The right-hand side is monotone in $\Delta^*$ by convexity of the absolute-value expectation in its argument. Differentiating implicitly in $\sigma_\eta^2$,
\[
\frac{\partial \Delta^*}{\partial \sigma_\eta^2} \;=\; \frac{-\,\partial^2 G/\partial S^{op}_1 \partial \sigma_\eta^2}{\partial^2 G / \partial (S^{op}_1)^2} \;>\; 0,
\]
since the numerator is positive (a larger forecast-error variance increases the marginal cost of an extra forward-committed MWh, because the realised $|\eta_t|$ contribution to imbalance does not net out across subperiods under per-period settlement) and the denominator is positive by convexity. This establishes \eqref{eq:shift}: the aggregator's quantity shift from forward to spot is strictly increasing in the forecast-error variance.

\section{Strategic bidder: optimal two-tranche ladder}\label{app:tranche}

We derive the strategic bidder's closed-form two-tranche ladder under uniform-price clearing, the FOCs in $(p^L, p^H, k^L)$, the equilibrium spread $\Delta p^* = \Delta_\mu/(2 b)$, and a stylisation note that flags the simplification absorbed into the boundary condition.

\subsection{Profit under uniform-price clearing}\label{app:tranche:profit}
The strategic bidder is the residual monopolist; there is no higher-priced fringe bid above its top tranche. Under uniform-price clearing (as in OMIE day-ahead and intraday auctions), the clearing price is therefore set by the marginal accepted tranche of the strategic bidder's own ladder: $p^L$ when only the low tranche clears, $p^H$ when both clear. For a delivery product at which the residual-demand intercept is $\bar A + \mu$ (with $\mu \sim \mathrm{Uniform}[-\Delta_\mu, +\Delta_\mu]$), the strategic bidder commits a ladder $\{(p^L, k^L), (p^H, k^H)\}$ with $p^L \le p^H$, $k^L + k^H = K$ \emph{before} $\mu$ realises. The cleared MWh at $\mu$ is
\[
q(\mu) = \begin{cases} k^L = q^L & \text{if } \mu \le \mu^*, \\ k^L + k^H = K & \text{if } \mu > \mu^*, \end{cases}
\qquad \mu^* \;\equiv\; k^L + b \, p^L - \bar A,
\]
where $b$ is the market thinness (either $b_{1t}$ or $b_{2t}$ depending on which market the ladder is on). The clearing price is uniformly $p^L$ when only the low tranche clears and uniformly $p^H$ when both clear --- the marginal accepted tranche price, applied to the entire accepted volume. Expected profit per delivery product is therefore
\begin{equation}
\mathbb{E}[\pi]
\;=\; \frac{1}{2\Delta_\mu}\!\left[\int_{-\Delta_\mu}^{\mu^*}\!(p^L - c)\,q^L\,d\mu \;+\; \int_{\mu^*}^{\Delta_\mu}\!(p^H - c)\,K\,d\mu\right],
\label{eq:Epi-uniform}
\end{equation}
with $(p^H - c)\,K$ in the second integrand reflecting that all $K$ units fetch $p^H$ when the high tranche is marginal.

\subsection{First-order conditions}\label{app:tranche:foc}
Take derivatives in $(p^L, p^H, k^L)$ holding $K = k^L + k^H$ fixed. Define $w^L = k^L/K$, $w^H = k^H/K$.

\paragraph{FOC in $p^H$.} $\frac{\partial \mathbb{E}[\pi]}{\partial p^H} = \frac{1}{2 \Delta_\mu} \int_{\mu^*}^{\Delta_\mu} K\, d\mu = \frac{K(\Delta_\mu - \mu^*)}{2\Delta_\mu} > 0$ unless $\mu^* = \Delta_\mu$; at an interior optimum the bidder pushes $p^H$ to the boundary of acceptance, with the binding condition $p^H = (\bar A + \Delta_\mu - K)/b$ at the highest $\mu$. Substituting $\mu^* = 0$ at the symmetric solution gives $p^{H*} = (\bar A + b c)/(2 b) + \Delta_\mu/(4 b)$.

\paragraph{FOC in $p^L$.}
\[
\frac{\partial \mathbb{E}[\pi]}{\partial p^L} \;=\; \frac{1}{2\Delta_\mu}\!\left[\int_{-\Delta_\mu}^{\mu^*}\!q^L \, d\mu \;+\; (p^L - c) q^L \cdot \frac{\partial \mu^*}{\partial p^L}\right] \;=\; \frac{1}{2\Delta_\mu}\!\left[(\mu^* + \Delta_\mu)\, q^L \;+\; (p^L - c) q^L\, b\right] \;=\; 0.
\]
Solving: $\mu^* + \Delta_\mu + b(p^L - c) = 0$, i.e.\ $\mu^* = -\Delta_\mu - b(p^L - c)$. Combined with the definition $\mu^* = k^L + b p^L - \bar A$, this gives $k^L = \bar A - b p^L + \mu^*$. At $\mu^* = 0$: $p^{L*} = (\bar A + b c)/(2 b) - \Delta_\mu/(4 b)$.

\paragraph{FOC in $k^L$.}
\[
\frac{\partial \mathbb{E}[\pi]}{\partial k^L} \;=\; \frac{p^L - c}{2 \Delta_\mu}(\mu^* + \Delta_\mu) \;-\; \frac{p^H - c}{2\Delta_\mu}(\Delta_\mu - \mu^*) \;=\; 0.
\]
At the symmetric optimum $\mu^* = 0$ this reduces to $p^L - c = p^H - c$ when $\Delta_\mu - \mu^* = \mu^* + \Delta_\mu$, but the asymmetric prices $p^L = p^{L*}$, $p^H = p^{H*}$ require $\mu^* = 0$ exactly. Substituting $\mu^* = 0$ in $k^L = \bar A - b p^{L*}$ and the closed form for $p^{L*}$,
\[
k^{L*} \;=\; \bar A - b\bigg(\frac{\bar A + bc}{2b} - \frac{\Delta_\mu}{4b}\bigg) \;=\; \frac{\bar A - b c}{2} + \frac{\Delta_\mu}{4}.
\]

\paragraph{Closed-form interior optimum.}
\begin{equation}
\boxed{\;p^{L*} = \frac{\bar A + b\, c}{2 b} - \frac{\Delta_\mu}{4 b},
\quad
p^{H*} = \frac{\bar A + b\, c}{2 b} + \frac{\Delta_\mu}{4 b},
\quad
\mu^{**} = 0,
\quad
k^{L*} = \frac{\bar A - b\, c}{2} + \frac{\Delta_\mu}{4}.\;}
\label{eq:tranche-spread}
\end{equation}

\paragraph{Comparative statics.} From \eqref{eq:tranche-spread}, the optimal pre-realisation tranche spread is
\begin{equation}
\Delta p^* \;=\; p^{H*} - p^{L*} \;=\; \frac{\Delta_\mu}{2 b}.
\label{eq:spread-cs}
\end{equation}
Running the same argument on the hourly aggregate $\tilde\mu = Q^{-1} \sum_t \mu_t$ (variance $\sigma_\mu^2/Q < \sigma_\mu^2$) gives a strictly smaller optimal spread, $\Delta p^*|_{\text{hourly}} = \tilde\Delta_\mu/(2 b)$ with $\tilde\Delta_\mu = \Delta_\mu/\sqrt{Q}$. The per-period ladder is therefore strictly richer than its hourly counterpart by a factor $\sqrt{Q}$.

\paragraph{Stylisation note.} The closed-form \eqref{eq:tranche-spread} requires that the capacity $K$ be tuned to the boundary condition $p^{H*}$ binds at the highest demand realisation $\mu = \Delta_\mu$; for general $K$ the FOC in $p^L$ also picks up a term $-bK(p^H-c)$ from the second integrand of \eqref{eq:Epi-uniform}, and the optimum shifts off the symmetric form. The closed form is therefore a stylised approximation, with the auction microfoundation absorbed into the boundary condition. Comparative statics in $\Delta_\mu$ and $b$ are robust to this stylisation: $\Delta p^*$ widens with the demand range and contracts with thinness.

Empirical measures of within-curve bid-shape dispersion (MW-weighted price SDs, effective tranche counts) are motivated by this widening but are not derived from the two-state reduction --- they require richer bid-curve structure than the two-tranche approximation supports.

\section{$(1,1)$ equilibrium}\label{app:state-60-60}

We solve the pre-reform game by backward induction: the Stage~2 best response (\S\ref{app:60-60:stage2}), then the four arbitrage regimes at Stage~1 (\S\ref{app:60-60:stage1}). The scarcity wedge $p_1 > p_2 > c$ survives in regimes (1), (3), and (4); only competitive full arbitrage closes it.

\subsection{Stage 2 best response}\label{app:60-60:stage2}
The strategic bidder's Stage~2 problem, taking $p_1$, $q_1$ and $s$ as given:
\[
\max_{p_2} (p_2 - c) q_2 \quad\text{s.t.}\quad q_2 = b_2(p_1 - p_2) + s - S^{op}_2.
\]
FOC: $b_2(p_1 - p_2) + s - S^{op}_2 + (p_2 - c)(-b_2) = 0$, giving
\begin{equation}
p_2^*(p_1, s) \;=\; \frac{p_1 + c}{2} + \frac{s - S^{op}_2}{2\, b_2}.
\label{eq:p2-60-60}
\end{equation}

\subsection{Stage 1 best response and equilibrium in $(p_1, s)$}\label{app:60-60:stage1}
Substituting \eqref{eq:p2-60-60} into the strategic bidder's total profit and using $q_1 = \bar A - b_1 p_1 - s - S^{op}_1$,
\[
\Pi(p_1, s) \;=\; (p_1 - c)(\bar A - b_1 p_1 - s - S^{op}_1) \;+\; (p_2^*(p_1, s) - c) q_2^*(p_1, s).
\]
\noindent The four arbitrage cases pin $s$ differently.

\paragraph{(R1) No arbitrage ($s = 0$).} The strategic bidder solves a Cournot problem in each market separately. The forward FOC gives $p_1^* = (\bar A - S^{op}_1 + b_1 c)/(2 b_1)$. Combining with \eqref{eq:p2-60-60} yields $p_1^* > p_2^* > c$ and $q_1, q_2 > 0$ (\citeauthor{ito_reguant_2016} Result~1).

\paragraph{(R2) Full arbitrage ($p_1 = p_2$).} Setting $p_1 = p_2$ in \eqref{eq:p2-60-60} gives $s = b_2(p_1 - c) + S^{op}_2$. Substituting back, total cleared MWh $q_1 + q_2 = \bar A - b_1 p_1 - S^{op}_1 - S^{op}_2 = \bar A - b_1 p_1 - \bar Y_h$ (using $S^{op}_1 + S^{op}_2 = \bar Y_h$). The strategic bidder maximises $(p_1 - c)(\bar A - b_1 p_1 - \bar Y_h)$, giving
\begin{equation}
p_1^{*,F} \;=\; \frac{\bar A - \bar Y_h + b_1 c}{2 b_1}, \qquad p_2^{*,F} \;=\; p_1^{*,F}.
\label{eq:p1-60-60-full}
\end{equation}
Note that $s > 0$ \emph{reduces} the strategic bidder's total output relative to no arbitrage (\citeauthor{ito_reguant_2016} Result~2): even competitive arbitrage does not restore competitive prices, because the bidder withholds more in response.

\paragraph{(R3) Limited arbitrage ($s \le K$ binding).} If $K$ is below the full-arbitrage level, $s = K$ binds and $p_1 - p_2 > 0$. From \eqref{eq:p2-60-60}, $p_1 - p_2 = (p_1 - c)/2 - (K - S^{op}_2)/(2 b_2)$; the wedge depends on $K$ and on $S^{op}_2$ (\citeauthor{ito_reguant_2016} Result~3). The smaller $K$, the larger the residual wedge.

\paragraph{(R4) Strategic arbitrage.} The arbitrageur maximises $\pi^A(s) = s(p_1 - p_2)$ subject to \eqref{eq:p2-60-60}:
\[
\pi^A(s) \;=\; s\!\left[\frac{p_1 - c}{2} - \frac{s - S^{op}_2}{2\, b_2}\right].
\]
FOC in $s$:
\[
\frac{p_1 - c}{2} - \frac{2 s - S^{op}_2}{2\, b_2} \;=\; 0
\quad\Longleftrightarrow\quad
s^{*,S}(p_1) \;=\; \frac{b_2 (p_1 - c)}{2} + \frac{S^{op}_2}{2}.
\]
Substituting into \eqref{eq:p2-60-60} (with $S^{op}_2 \to 0$ for clarity in the wedge derivation),
\begin{equation}
p_2^{*,S} \;=\; \frac{p_1 + c}{2} + \frac{b_2(p_1 - c)/2}{2\, b_2} \;=\; \frac{p_1 + c}{2} + \frac{p_1 - c}{4} \;=\; \frac{3 p_1 + c}{4}.
\label{eq:p2-strategic-derived}
\end{equation}
The forward-spot wedge is $p_1 - p_2^{*,S} = (p_1 - c)/4$, matching \eqref{eq:strategic-arb} in the main body. (\citeauthor{ito_reguant_2016} Result~4.) The Stage~1 forward price $p_1^{*,S}$ is then the maximiser of the strategic bidder's compound profit; algebra in the standard quadratic gives $p_1^{*,S} - c = 8 (\bar A - b_1 c - \bar Y_h)/(16 b_1 - b_2)$ at $S^{op}_2 = 0$, $S^{op}_1 = \bar Y_h$. The wedge $(p_1^{*,S} - c)/4 > 0$ persists at the equilibrium price.

\section{$(1,Q)$ equilibrium}\label{app:state-60-15}

Forward hourly, spot per-period. We derive the per-product spot best response, work through full and strategic arbitrage to obtain the within-block dispersion equations, and finally compute the within-hour wedge SD in this asymmetric state.

The strategic bidder's Stage~2 problem now has $Q$ spot products. Directly from the spot residual demand $q_{2t} = b_{2t}(p_1 - p_{2t}) + s_t - S^{op}_{2t}$ in \eqref{eq:rd} and the FOC $q_{2t} = b_{2t}(p_{2t}-c)$, the strategic per-period spot best response is
\begin{equation}
p_{2t}^*(p_1, s_t) \;=\; \frac{p_1 + c}{2} \;+\; \frac{s_t - S^{op}_{2t}}{2\, b_{2t}}.
\label{eq:p2-60-15-foc}
\end{equation}
The cross-product variation in $p_{2t}^*$ in this state comes purely from per-subperiod variation in $b_{2t}$, $s_t$, $S^{op}_{2t}$ --- the forward price $p_1$ is hourly and shared, and $\mu_t$ does not enter the spot demand directly (it acts only through the forward-cleared volume, which is uniform per subperiod when forward is hourly).

\subsection{Full arbitrage: block parity and within-block dispersion}\label{app:60-15:full}
Block parity requires $p_1 = (1/Q) \sum_t p_{2t}$. Substituting \eqref{eq:p2-60-15-foc},
\[
p_1 \;=\; \frac{p_1 + c}{2} + \frac{1}{2 Q}\!\sum_t\!\frac{s_t - S^{op}_{2t}}{b_{2t}}
\quad\Longleftrightarrow\quad
p_1 - c \;=\; \frac{1}{Q}\!\sum_t\!\frac{s_t - S^{op}_{2t}}{b_{2t}}.
\]
Setting $S^{op}_{2t} \equiv 0$ and uniform arbitrage $s_t \equiv \bar s$ to isolate the thinness mechanism, and $Q = 2$ for concreteness,
\[
p_1 - c \;=\; \frac{\bar s}{2}\!\left(\frac{1}{b_{21}} + \frac{1}{b_{22}}\right) \;=\; \frac{\bar s}{2}\cdot\frac{b_{21}+b_{22}}{b_{21} b_{22}}
\quad\Longleftrightarrow\quad
\bar s \;=\; \frac{2(p_1 - c)\, b_{21} b_{22}}{b_{21} + b_{22}}.
\]
The within-block dispersion follows from differencing $p_{2t}^*$ across $t$:
\[
p_{21} - p_{22} \;=\; \bigg(\frac{\bar s}{2 b_{21}} - \frac{\bar s}{2 b_{22}}\bigg) \;=\; \frac{\bar s}{2}\cdot\frac{b_{22} - b_{21}}{b_{21}\, b_{22}}.
\]
Substituting $\bar s$ from above,
\begin{equation}
\boxed{\;p_{21} - p_{22} \;=\; (p_1 - c)\cdot \frac{b_{22} - b_{21}}{b_{21} + b_{22}}\;}\qquad\text{(full arbitrage)}.
\label{eq:within-block-full}
\end{equation}

\subsection{Strategic arbitrage: block wedge and attenuated dispersion}\label{app:60-15:strategic}
The strategic arbitrageur internalises price impact and stops at
\[
p_1 - (1/Q)\sum_t p_{2t} \;=\; \frac{p_1 - c}{4}.
\]
The block-wedge condition is the average of \eqref{eq:strategic-arb} across the $Q$ matched spot products. Repeating the algebra of \S\ref{app:60-15:full} with this modified parity, the arbitrage position $\bar s$ falls to exactly half of the full-arbitrage value:
\[
\bar s^{S} \;=\; \frac{(p_1 - c)\, b_{21} b_{22}}{b_{21} + b_{22}} \;=\; \frac{\bar s^{F}}{2}.
\]
The within-block dispersion correspondingly becomes
\begin{equation}
\boxed{\;p_{21} - p_{22} \;=\; \frac{p_1 - c}{2}\cdot\frac{b_{22} - b_{21}}{b_{21} + b_{22}}\;}\qquad\text{(strategic arbitrage)}.
\label{eq:within-block-strategic}
\end{equation}
The dispersion is attenuated by a factor of two relative to the full-arbitrage case, reflecting the smaller arbitrage volume traded under strategic incentives. The mechanism is dampened but not eliminated.

\subsection{Two-tranche ladder on the spot side}\label{app:60-15:ladder}
The strategic bidder posts per-product spot ladders. Applying \S\ref{app:tranche} with $b \to b_{2t}$ gives ladder spread $\Delta p^*_{2t} = \Delta_\mu/(2 b_{2t})$, strictly larger than the spread under hourly spot (which uses $\tilde\Delta_\mu = \Delta_\mu / \sqrt{Q}$). This is the spot-side richer-ladder mechanism reported in the main body, eq.~\eqref{eq:tranche-id}.

\subsection{Wedge SD in the asymmetric state}\label{app:60-15:wedgesd}
The per-subperiod wedge $w_t = p_1 - p_{2t}^*$ inherits the per-period variation in $p_{2t}^*$; $p_1$ is constant across $t$. From \eqref{eq:p2-60-15-foc}, $p_{2t}^*$ varies across $t$ only through per-subperiod thinness $b_{2t}$ (and the arbitrage and aggregator terms when those vary by $t$). For two matched products ($Q=2$) the population SD of $w_t$ equals half the within-block price difference, so the full and strategic arbitrage benchmarks give
\begin{equation}
\mathrm{sd}_t^{F}(w_t) \;=\; \frac{(p_1 - c)\,|b_{22} - b_{21}|}{2(b_{21} + b_{22})}, \qquad
\mathrm{sd}_t^{S}(w_t) \;=\; \frac{(p_1 - c)\,|b_{22} - b_{21}|}{4(b_{21} + b_{22})},
\label{eq:wedge-sd-asym-full-strat}
\end{equation}
the strategic value exactly half of the full-arbitrage value, reflecting the smaller arbitrage volume traded under strategic incentives. Under limited arbitrage with per-product capacity $s_t = K/Q$ binding, the wedge becomes $w_t = (p_1-c)/2 - K/(2Q b_{2t})$, and
\begin{equation}
\mathrm{sd}_t^{L}(w_t) \;=\; \frac{K\,|b_{22} - b_{21}|}{8\, b_{21}\, b_{22}} \qquad (Q=2,\; s_t = K/2),
\label{eq:wedge-sd-asym-limited}
\end{equation}
which scales with $K$ rather than with $(p_1-c)$ but inherits the same $|b_{22}-b_{21}|$ thinness asymmetry. All three benchmarks scale with per-subperiod thinness asymmetry rather than with $\sigma_\mu$ --- consistent with the main-body characterisation in \S\ref{sec:id15}(M3).

\section{$(Q,Q)$ equilibrium}\label{app:state-15-15}

Both markets per-period. We derive the strategic forward best response per product, work through full, strategic, and (the empirically relevant) limited-arbitrage regimes, and close with the Cournot scale-up of cleared MWh and the Stage-1 option-value mechanism behind the cross-market substitution of bid shapes.

The strategic per-period forward best response from \eqref{eq:rd} with $\bar A_t = \bar A + \mu_t$:
\begin{equation}
p_{1t}^* \;=\; \frac{\bar A + \mu_t - S^{op}_{1t} + b_{1t}\, c}{2\, b_{1t}}.
\label{eq:p1t-15-15}
\end{equation}

\subsection{Full arbitrage: product parity}\label{app:15-15:full}
Per-product parity $p_{1t} = p_{2t}$ pins $s_t$ at $s_t = b_{2t}(p_{1t} - c) + S^{op}_{2t}$ (same structure as the $(1,1)$ full-arbitrage condition, but per product). The wedge SD is zero by construction.

\subsection{Strategic arbitrage: product wedge and SD}\label{app:15-15:strategic}
Per-product strategic arbitrage gives \eqref{eq:strategic-arb} applied per product: $p_{2t} = (3 p_{1t} + c)/4$. The per-product wedge is $p_{1t} - p_{2t} = (p_{1t} - c)/4$. Cross-period SD:
\[
\mathrm{sd}_t(p_{1t} - p_{2t}) \;=\; \tfrac{1}{4}\,\mathrm{sd}_t(p_{1t}^*).
\]
From \eqref{eq:p1t-15-15} with $S^{op}_{1t}$ approximately matched to $\bar y_t$ (so its variation absorbs structural production variation; the remaining driver is $\mu_t$):
\[
\mathrm{sd}_t(p_{1t}^*) \;=\; \mathrm{sd}_t\!\bigg(\frac{\mu_t}{2 b_1}\bigg) \;=\; \frac{\sigma_\mu}{2\, b_1}.
\]
Therefore
\begin{equation}
\boxed{\;\mathrm{sd}_t(p_{1t} - p_{2t}) \;=\; \frac{1}{4}\cdot\frac{\sigma_\mu}{2\, b_1} \;=\; \frac{\sigma_\mu}{8\, b_1}.\;}
\label{eq:wedge-sd-15-15}
\end{equation}
This is strictly smaller than the asymmetric-state wedge SD $\sigma_\mu/(2 b_2)$ whenever $b_2 \le 4 b_1$, a condition that holds under the \citeauthor{ito_reguant_2016} assumption $b_2 \le b_1$.

\subsection{Two-tranche ladder on the forward side}\label{app:15-15:ladder}
Applying \S\ref{app:tranche} with $b \to b_{1t}$ gives forward-side per-product ladder spread $\Delta p^*_{1t} = \Delta_\mu/(2 b_{1t})$, mirroring the spot-side result of \S\ref{app:60-15:ladder}.

\subsection{Limited arbitrage at the symmetric state (empirical regime)}\label{app:15-15:limited}

The strategic-arbitrage wedge \eqref{eq:wedge-sd-15-15} assumes the arbitrageur can build any per-product position. With $Q$ matched products per clock-hour the arbitrageur's capacity $K$ is spread over $Q$ positions, so the effective per-product capacity is $K_t \approx K/Q$. When this capacity binds --- which is the empirically relevant case in electricity (see \S\ref{sec:speaks} ``Arbitrage regime is empirical, not theoretical'') --- the per-product equilibrium is \citeauthor{ito_reguant_2016} Result~3 (limited arbitrage) applied product by product.

\paragraph{Strategic spot best response under binding per-product capacity.} Setting $s_t = K_t$ in the strategic Stage~2 FOC \eqref{eq:p2-60-15-foc}:
\[
p_{2t}^* \;=\; \frac{p_{1t}^* + c}{2} + \frac{K_t - S^{op}_{2t}}{2\, b_{2t}}.
\]
The forward per-product cleared price $p_{1t}^*$ depends on $\mu_t$ through the forward demand $\bar A_t = \bar A + \mu_t$; from \eqref{eq:p1t-15-15} (corrected to include the arbitrage position), $p_{1t}^* = (\bar A + \mu_t - S^{op}_{1t} - s_t + b_{1t} c)/(2 b_{1t})$ so $\mathrm{sd}_t(p_{1t}^*) = \sigma_\mu/(2 b_1)$ (under uniform $s_t = K_t$ and $S^{op}_{1t}$ matched to structural production). The per-product wedge is
\[
p_{1t}^* - p_{2t}^* \;=\; \frac{p_{1t}^* - c}{2} \;-\; \frac{K_t - S^{op}_{2t}}{2\, b_{2t}}.
\]

\paragraph{Wedge SD.} Taking the cross-product SD:
\[
\mathrm{sd}_t(p_{1t}^* - p_{2t}^*) \;=\; \tfrac{1}{2}\,\mathrm{sd}_t(p_{1t}^*) \;=\; \frac{\sigma_\mu}{4\, b_1}.
\]
Compare to the strategic-arbitrage closed-form \eqref{eq:wedge-sd-15-15}: under unlimited strategic arbitrage the wedge is $(p_{1t}-c)/4$ giving wedge SD $\sigma_\mu/(8\, b_1)$; under limited arbitrage the wedge is $(p_{1t}-c)/2$ giving wedge SD $\sigma_\mu/(4\, b_1)$, exactly twice the strategic-arbitrage value. The wedge SD does \emph{not} collapse fully even at the symmetric state. If per-period thinness $b_{2t}$ also varies across $t$ (as in the asymmetric state) the wedge inherits an additional thinness-asymmetry component on top of $\sigma_\mu/(4 b_1)$; the empirically observed sustaining of wedge SD at granular-state magnitudes is consistent with this combined effect.

\paragraph{Per-product capacity binding.} For the limited regime to bind, $K_t$ must be smaller than the strategic-arbitrage equilibrium $s_t^{\text{strat}} = b_{2t}(p_{1t}-c)/2$. With $K_t \approx K/Q$ and $s_t^{\text{strat}}$ proportional to a per-product mark-up, binding requires $K \le Q\cdot \bar s^{\text{strat}}$ where $\bar s^{\text{strat}}$ is the average per-product strategic position. Empirically, hourly arbitrage capacity is set by participation rules and balance-sheet limits that do not scale with $Q$ when granularity reforms multiply products; the constraint binds at the granular state by construction.

\subsection{Cournot cleared-MWh under per-period thinness}\label{app:15-15:scaleup}

The strategic best response \eqref{eq:p1t-15-15} gives per-period cleared quantity $q_{1t}^* = (\bar A_t - b_{1t} c - S^{op}_{1t} - s_t)/2$. Summing across $t$ and taking expectation over $\mu_t$ (so $\mathbb{E}[\sum_t \bar A_t] = Q\bar A$):
\[
\mathbb{E}\!\bigg[\sum_t q_{1t}^*\bigg] \;=\; \frac{Q \bar A - c \sum_t b_{1t} - \sum_t S^{op}_{1t} - \sum_t s_t}{2}.
\]
The hourly counterpart is $Q\cdot q_1^*|_{(1,1)} = Q\cdot(\bar A - b_1 c - S^{op}_1 - s)/2$. Comparing (with the aggregator and arbitrage totals comparable across configurations):
\begin{equation}
\mathbb{E}\!\bigg[\sum_{t=1}^{Q} q_{1t}^*\bigg]\bigg|_{(Q,Q)} \;-\; Q\cdot q_1^*\big|_{(1,1)} \;=\; \frac{c}{2}\,\bigg(Q\, b_1 - \sum_{t=1}^{Q} b_{1t}\bigg).
\label{eq:scale-up}
\end{equation}
Positive whenever $\sum_{t=1}^{Q} b_{1t} < Q\,b_1$ --- i.e.\ when per-period thinness is on average smaller than the hourly aggregate elasticity, as expected (each 15-min product is thinner than the hourly aggregate it replaces). This is the Cournot scale-up of main-body \S\ref{sec:da15}(iv), obtained without any imported mechanism.

\subsection{Stage-1 option value (cross-market substitution)}\label{app:15-15:option}

Stage-1 expected profit including Stage~2 is, at second order in $\mu_t$,
\[
\mathbb{E}[\Pi(p_1)] \;=\; \pi_1(p_1) \;+\; \mathbb{E}_{\mu}[\pi_2(p_1, \mu_t)] \;=\; \pi_1(p_1) \;+\; \pi_2(p_1, 0) \;+\; \tfrac{1}{2}\sigma_\mu^2 \cdot \frac{\partial^2 \pi_2}{\partial \mu_t^2}\bigg|_{\mu_t = 0}\quad\text{(quadratic-in-}\mu_t\text{; no higher-order remainder)}.
\]
Per-period spot Cournot makes $\pi_2$ convex in $\mu_t$ (since the per-product cleared quantity scales with $\mu_t/(2 b_{2t})$ and the per-product mark-up scales with $\mu_t/(2 b_{2t})$ via \eqref{eq:p2-60-15-foc}, the product is quadratic in $\mu_t$). The second derivative is positive: the variance correction \emph{raises} Stage-1 expected profit. The Stage-1 FOC in $p_1$ shifts accordingly: the bidder optimises against a more valuable Stage-2 option, widening the forward ladder to capture both sides of the variation. Sign is unambiguous; magnitude depends on $\sigma_\mu^2/(b_{1t} b_{2t})$. Symmetrically when the forward is per-period: the within-hour discrimination is absorbed at Stage~1 directly, so the Stage-2 option value drops and the spot ladder tightens. This is the model-internal account of the cross-market substitution pattern of empirical \S 4.C.3.

\section{Balancing discrepancy}\label{app:balancing}

We collect the balancing-stage discrepancy under each state, summing demand-side unmet load ($\mu_t$ that no auction absorbed) and aggregator production-side unmet schedule ($\eta_t$ that no per-period spot absorbed). These two channels are what the imbalance penalty $\gamma$ prices in welfare (\S\ref{app:welfare}).

The balancing-stage discrepancy is the difference between realised demand plus aggregator under-/over-production and the auction-scheduled supply, summed per subperiod. With $\bar A_t = \bar A + \mu_t$ and aggregator realisation $y_t = \bar y_t + \eta_t$,
\begin{equation}
\Delta_t \;=\; (\bar A_t - q^{\text{cleared}}_t) \;+\; (y_t - S^{op}_t) \;+\; \varepsilon_t,
\label{eq:balancing-disc}
\end{equation}
where $S^{op}_t = S^{op}_{1t} + S^{op}_{2t}$ is the aggregator's total scheduled MWh at subperiod $t$ and $\varepsilon_t$ is residual noise unrelated to $\mu_t$ and $\eta_t$.

\emph{State $(1,1)$.} Forward and spot are both hourly: the forward and spot schedules clear uniformly per subperiod, so per-subperiod demand variation $\mu_t$ is unabsorbed; the aggregator commits to a Stage-1 uniform per-subperiod supply and the realisation shock $\eta_t$ loads onto its schedule deviation. Under the structural assumption $\bar y_t \equiv \bar Y_h / Q$ (no within-hour structural production variation), $\Delta_t = \mu_t + \eta_t + \varepsilon_t$.

\emph{State $(1,Q)$.} The spot becomes per-period: the aggregator re-schedules at Stage~2 using its updated forecast $\bar y_t + \eta_t$, absorbing the $\eta_t$ shock into its spot schedule (\S\ref{app:aggregator:60-15}). The forward remains hourly so $\mu_t$ is still unabsorbed: $\Delta_t = \mu_t + \varepsilon_t$.

\emph{State $(Q,Q)$.} The forward also becomes per-period: per-product forward clearing $q_{1t}$ responds to $\bar A_t = \bar A + \mu_t$ through \eqref{eq:p1t-15-15}, absorbing $\mu_t$. The aggregator's $\eta_t$ remains absorbed at spot per-period. The balancing discrepancy reduces to $\Delta_t = \varepsilon_t$.

These reductions are mechanical in the contracting grid and hold under any arbitrage regime.

\section{Note on the alternative sequence}\label{app:alt-sequence}

The mirror path $(1,1) \to (Q,1) \to (Q,Q)$ swaps the order of the two auction reforms. After relabelling, the same per-state algebra of \S\ref{app:state-60-60}--\ref{app:state-15-15} applies. The end-state $(Q,Q)$ is identical. The intermediate state $(Q,1)$ produces a within-hour \emph{DA} dispersion driven by per-subperiod $b_{1t}$ asymmetry (mirror of \eqref{eq:within-block-full}--\eqref{eq:within-block-strategic} but with roles of $b_{1t}$ and $b_{2t}$ swapped); the aggregator's quantity-shift mechanism activates already at the forward reform; the balancing-discrepancy gain $\Delta_t : \mu_t + \varepsilon_t \to \varepsilon_t$ realises at the forward reform under either path. The peak wedge SD ($\sigma_\mu / (2\, b_2)$ in the original path, $\sigma_\mu/(2\, b_1)$ in the mirror) and the final wedge SD ($\sigma_\mu/(8\, b_1)$ in both) are sequence-invariant up to the relabelling of the dominant elasticity.

\section{Welfare derivations}\label{app:welfare}

We collect the closed-form welfare components by state and arbitrage regime under uniform $\mu_t \sim [-\Delta_\mu, \Delta_\mu]$ and $\eta_t \sim [-\Delta_\eta, \Delta_\eta]$ (residual noise $\varepsilon_t$ negligible). The aggregator's penalty enters $\pi^{Agg,\text{net}}$ through $\eta_t$ and the system imbalance cost $C^I_{\text{sys}}$ enters through $\mu_t$; the two channels are disjoint.

\subsection{Components}\label{app:welfare:components}

Per clock-hour:
\begin{equation}
W \;=\; CS \;+\; \pi^{SB} \;+\; \pi^{Agg,\text{net}} \;+\; \pi^{Arb} \;-\; C^I_{\text{sys}}.
\label{eq:welfare-app}
\end{equation}
The aggregator's net profit $\pi^{Agg,\text{net}}$ is its revenue $\sum_t (p_1^{(t)} S^{op}_{1t} + p_{2t} S^{op}_{2t})$ minus its internalised quadratic penalty $\gamma\, E[\mathcal{I}^{Agg}] = \gamma \sum_t E[(\eta_t^{\text{unhedged}})^2]$ (\S\ref{app:aggregator}). The system imbalance cost is $C^I_{\text{sys}} = \gamma \sum_t E[\Delta_t^2]$ with $\Delta_t$ from \S\ref{app:balancing}. $\pi^{Arb} = s(p_1 - p_2)$ is the arbitrageur's revenue ($\pi^{Arb} = 0$ under no arbitrage and full arbitrage). For linear residual demand $q = \bar A - b p$, consumer surplus at $p^*$ is
\begin{equation}
CS \;=\; \frac{(\bar A - b p^*)^2}{2 b}.
\label{eq:cs-formula}
\end{equation}

\subsection{Imbalance cost decomposition (closed form, uniform shocks)}\label{app:welfare:imbalance}

Under $\mu_t \sim U[-\Delta_\mu, \Delta_\mu]$ and $\eta_t \sim U[-\Delta_\eta, \Delta_\eta]$ independent (and $\varepsilon_t \approx 0$): $E[\mu_t^2] = \Delta_\mu^2/3$, $E[\eta_t^2] = \Delta_\eta^2/3$. The aggregator's unhedged shock (\S\ref{app:aggregator:hourly}) is $\eta_t$ when spot is hourly and $0$ once spot is per-period; the system discrepancy (\S\ref{app:balancing}) is $\mu_t$ when forward is hourly and $0$ once forward is per-period. The quadratic-penalty costs decompose as
\begin{equation}
\renewcommand{\arraystretch}{1.2}
\begin{array}{l c c c}
& (1,1) & (1,Q) & (Q,Q) \\ \hline
\gamma\, E[\mathcal{I}^{Agg}]\ \text{(in }\pi^{Agg,\text{net}}\text{)} & \gamma Q \Delta_\eta^2 / 3 & 0 & 0 \\
C^I_{\text{sys}}\ \text{(separate)} & \gamma Q \Delta_\mu^2 / 3 & \gamma Q \Delta_\mu^2 / 3 & 0 \\
\text{Total imbalance cost} & \gamma Q(\Delta_\mu^2 + \Delta_\eta^2)/3 & \gamma Q \Delta_\mu^2 / 3 & 0
\end{array}
\label{eq:imb-table}
\end{equation}
The two channels are disjoint ($\eta_t$ only enters $\pi^{Agg,\text{net}}$, $\mu_t$ only enters $C^I_{\text{sys}}$). The two reform steps save $\gamma Q \Delta_\eta^2/3$ (spot, absorbed into $\Delta \pi^{Agg,\text{net}}$) and $\gamma Q \Delta_\mu^2/3$ (forward, absorbed into $-\Delta C^I_{\text{sys}}$).

\subsection{Welfare at state $(1,1)$ by arbitrage regime}\label{app:welfare:1-1}

\paragraph{No arbitrage (R1).} From \S\ref{app:60-60:stage1}: $p_1^{NA} = (\bar A + b_1 c - S^{op}_1)/(2 b_1)$, $q_1^{NA} = (\bar A - b_1 c - S^{op}_1)/2$, $\pi^{SB,NA}_1 = (\bar A - b_1 c - S^{op}_1)^2/(4 b_1)$. Consumer surplus at the forward stage is, by \eqref{eq:cs-formula}, $CS^{NA}_1 = (\bar A - b_1 p_1^{NA})^2/(2 b_1) = (q_1^{NA} + S^{op}_1)^2/(2 b_1)$. The spot stage adds analogous terms via the spot best response \eqref{eq:p2-60-60}. Total welfare:
\[
W^{NA}_{(1,1)} \;=\; CS^{NA}_1 + CS^{NA}_2 + \pi^{SB,NA}_1 + \pi^{SB,NA}_2 + \pi^{Agg,\text{net},NA} - \gamma Q \Delta_\mu^2/3,
\]
with $\pi^{Agg,\text{net},NA}$ including the $-\gamma Q \Delta_\eta^2/3$ internalised penalty.

\paragraph{Full arbitrage (R2).} At $p_1 = p_2 = p_1^{F} = (\bar A - \bar Y_h + b_1 c)/(2 b_1)$ from \eqref{eq:p1-60-60-full}, strategic-bidder profit is $\pi^{SB,F} = (p_1^F - c)(\bar A - b_1 p_1^F - \bar Y_h) = (\bar A - b_1 c - \bar Y_h)^2/(4 b_1)$. The cleared MWh stays below the competitive benchmark (\citeauthor{ito_reguant_2016} Result~2): with monopoly market power, even competitive arbitrage cannot restore the first-best.

\paragraph{Strategic arbitrage (R4).} The arbitrageur extracts $\pi^{Arb,S} = s^{*,S}(p_1 - p_2)|_{s = s^{*,S}} = b_2(p_1 - c)^2/4$ (\S\ref{app:60-60:stage1}), a pure transfer that subtracts from the strategic bidder's rent. Welfare lies strictly between $W^{NA}$ and $W^{F}$.

\paragraph{Limited arbitrage (R3, empirical).} With $s = K < b_2(p_1 - c)/2$ binding, the arbitrageur extracts $\pi^{Arb,L} = K\, [(p_1 - c)/2 - (K - S^{op}_2)/(2 b_2)]$. Welfare is closer to $W^{NA}$ when $K$ is small, and converges to $W^{S}$ as $K$ grows to $b_2(p_1-c)/2$.

\subsection{Welfare at $(1,Q)$ and $(Q,Q)$ --- delta from $(1,1)$}\label{app:welfare:reform}

The reform's per-step welfare change decomposes into closed-form imbalance, ladder-rent, and (at $(Q,Q)$) Cournot scale-up terms.

\paragraph{Spot reform, $(1,1)\to(1,Q)$: aggregator hedging.} The aggregator re-schedules per-subperiod at spot \eqref{eq:agg-perperiod-id}; $E[\mathcal{I}^{Agg}]$ drops from $Q\,\Delta_\eta^2/3$ to $0$, so $\Delta \pi^{Agg,\text{net}} = +\gamma Q\,\Delta_\eta^2/3$.

\paragraph{Spot-side ladder rent.} The strategic bidder's spot-side two-tranche ladder \eqref{eq:tranche-id} replaces a single-price Cournot with a $\Delta_\mu$-discriminated ladder. By \S\ref{app:tranche} closed form (substituting $b = b_{2t}$):
\begin{equation}
\Delta \pi^{SB,2T}_{\text{spot}} \;=\; \frac{\Delta_\mu\,\big[4(\bar A - b_{2t} c) + \Delta_\mu\big]}{16\, b_{2t}}.
\label{eq:ladder-rent-spot}
\end{equation}
On the inframarginal volume $k^{L*} = (\bar A - b_{2t} c)/2 + \Delta_\mu/4$, this is a CS-to-$\pi^{SB}$ transfer; on the marginal extension of $\Delta_\mu/2$ MWh, this is new surplus. The net welfare gain is $\Delta_\mu(\bar A - b_{2t}c)/(4 b_{2t}) + \Delta_\mu^2/(16 b_{2t})$ from the extension --- a fraction of \eqref{eq:ladder-rent-spot}.

\paragraph{Forward reform, $(1,Q)\to(Q,Q)$: system imbalance.} Per-product forward absorbs $\mu_t$ via \eqref{eq:p1t-15-15}, so $C^I_{\text{sys}}$ drops from $\gamma Q\,\Delta_\mu^2/3$ to $0$:
\[
\Delta(-C^I_{\text{sys}}) \;=\; +\gamma Q\,\Delta_\mu^2/3.
\]

\paragraph{Forward-side ladder rent.} The same closed form as \eqref{eq:ladder-rent-spot} applies on the forward side with $b = b_1$:
\begin{equation}
\Delta \pi^{SB,2T}_{\text{forward}} \;=\; \frac{\Delta_\mu\,\big[4(\bar A - b_1 c) + \Delta_\mu\big]}{16\, b_1}.
\label{eq:ladder-rent-forward}
\end{equation}

\paragraph{Cournot scale-up at $(Q,Q)$.} Total cleared MWh per clock-hour rises by $(c/2)(Q b_1 - \sum_t b_{1t})$ via \eqref{eq:scale-up}. At the mark-up $p_1 - c = (\bar A - b_1 c - S^{op}_1)/(2 b_1)$,
\begin{equation}
\Delta W^{\text{scale}} \;=\; \frac{c\,(\bar A - b_1 c - S^{op}_1)}{4\, b_1}\,\Big(Q b_1 - \sum_t b_{1t}\Big).
\label{eq:scaleup-welfare-app}
\end{equation}

\paragraph{Limited-arbitrage gain at $(1,Q)$ and $(Q,Q)$.} Per-product capacity $K/Q$ binds, leaving a per-product wedge $(p_{1t}-c)/2$ instead of the unlimited-strategic $(p_{1t}-c)/4$ (\S\ref{app:15-15:limited}). The strategic bidder's spot markup over marginal cost \emph{falls} from $3(p_1-c)/4$ to $(p_1-c)/2$, so the per-product deadweight loss falls by $5 b_2(p_1-c)^2/32$ (\citet{ito_reguant_2016} Result~2: less arbitrage means less strategic withholding). This is the welfare counterpart of the empirical asymmetric-persistence wedge-SD finding --- not a welfare cost.

\subsection{Net welfare deltas (closed form)}\label{app:welfare:deltas}

Combining the channels for each reform step (writing $a_m \equiv \bar A - b_m c$ for $m \in \{1, 2\}$):
\begin{align}
\Delta W^{\text{spot}} \;&=\; \frac{\gamma Q \Delta_\eta^2}{3} \;+\; \frac{5\, b_2 (p_1-c)^2}{32} \;+\; \frac{\Delta_\mu\,[4 a_2 + \Delta_\mu]}{16\, b_2}, \label{eq:DeltaW-spot-app}\\
\Delta W^{\text{forward}} \;&=\; \frac{\gamma Q \Delta_\mu^2}{3} \;+\; \frac{c\, a_1}{4\, b_1}\,\Big(Q b_1 - \!\sum_t b_{1t}\Big) \;+\; \frac{\Delta_\mu\,[4 a_1 + \Delta_\mu]}{16\, b_1}, \label{eq:DeltaW-fwd-app}
\end{align}
where the second term of \eqref{eq:DeltaW-spot-app} is the per-product limited-arbitrage gain at the regime change, and the third terms in both lines are the bid-ladder rent extensions (mostly distributional).

\paragraph{Sign and dominance.} Both deltas are unconditionally positive --- the model predicts welfare improvement at each reform step, with imbalance saving the dominant channel given that balancing prices in Spanish electricity markets exceed day-ahead clearing prices by a factor of $2$--$10$.

\end{document}
